\chapter{Tingling Numbness}

\section*{Definition}
Abnormal sensation termed paresthesia tingling prickling burning loss normal skin sensation affecting extremities face trunk

\section*{When to See a Doctor}
See your doctor if:
\begin{itemize}
  \item Tingling or numbness is accompanied by weakness, speech difficulty, or coordination problems
\end{itemize}

\section*{How It Develops}
Results nerve compression ischemia demyelination peripheral neuropathy central nervous system lesions stroke TIA vitamin deficiency multiple sclerosis carpal tunnel cervical radiculopathy disrupted sensory pathway abnormal firing patterns nerve signaling

Diabetic neuropathy presents burning shooting pains tingling numbness loss coordination affecting feet legs sometimes hands arms more common with poor diabetes control

Peripheral neuropathy typically begins finger toe tips slowly progresses along limbs causes burning shooting pains tingling numbness

Nutritional neuropathy associated vitamin deficiencies causes burning tingling sensory symptoms plus sore red tongue fatigue anemia signs

\section*{Causes}
\begin{itemize}
  \item Carpal tunnel wrist numbness nocturnal predominant
  \item Peripheral neuropathy diabetes alcohol vitamin B12 deficiency
  \item Cervical radiculopathy neck arm dermatomal distribution
  \item Transient ischemic attack stroke sudden focal deficits
  \item Multiple sclerosis young adults relapsing remitting episodic
  \item Vitamin B12 deficiency pernicious anemia folate deficiency
  \item Alcohol toxic neuropathy chronic heavy consumption
\end{itemize}

\section*{Related Symptoms}
\begin{itemize}
  \item Weakness muscle twitching burning sensation allodynia distribution patterns dermatome nerve territory localized affected body part
\end{itemize}

\section*{Medical Evaluation}
\begin{itemize}
  \item Detailed neurological history includes onset pattern, progression, triggers, associated symptoms, and functional impact.
  \item Comprehensive neurological examination assesses mental status, cranial nerves, motor/sensory function, reflexes, coordination, and gait.
  \item Neuroimaging (CT, MRI) is often indicated to evaluate for structural, vascular, or demyelinating causes.
  \item Specialized testing (EEG for seizures, nerve conduction studies for neuropathy, lumbar puncture for meningitis) is guided by clinical suspicion.
\end{itemize}

\section*{Treatment Options}
\begin{itemize}
  \item Treat the underlying cause, such as diabetes, vitamin deficiency, medication toxicity, nerve compression, or inflammatory neuropathy.
  \item Neuropathic pain medicines may help painful symptoms, but they do not replace evaluation for progressive weakness, spinal-cord compression, or stroke.
  \item Physical or occupational therapy and neurology referral may be appropriate for persistent functional impairment or an uncertain diagnosis.
\end{itemize}

\section*{Self-Care and Home Management}
\begin{itemize}
  \item Protect the affected area from injury — you may not feel pain, heat, or cold normally.
  \item Avoid prolonged pressure on numb areas (don't sit cross-legged or lean on elbows for too long).
  \item If tingling affects your hands, be careful with sharp objects and hot surfaces.
  \item Check your feet daily for cuts or blisters if numbness affects your legs, especially if you have diabetes.
  \item Avoid driving if the numbness affects your ability to feel pedals or controls.
  \item Report persistent or progressive numbness to your doctor for evaluation.
\end{itemize}

\section*{References}
\begin{thebibliography}{99}
\bibitem{tingling_numbness:periphneuro1} England JD, Gronseth GS, Franklin G, et al. ``Practice parameter: evaluation of distal symmetric polyneuropathy.'' Neurology. 2009;72(2):185-192.
\bibitem{tingling_numbness:periphneuro2} Callaghan BC, Cheng HT, Stables CL, et al. ``Diabetic neuropathy: a review.'' JAMA. 2012;308(22):2371-2381.
\bibitem{tingling_numbness:periphneuro3} Watson JC, Dyck PJB. ``Peripheral neuropathy: a practical approach.'' Mayo Clin Proc. 2015;90(7):940-951.
\bibitem{tingling_numbness:periphneuro4} Azhary H, Farooq MU, Bhanushali M, et al. ``Peripheral neuropathy: differential diagnosis and management.'' Am Fam Physician. 2010;81(7):887-892.
\bibitem{tingling_numbness:periphneuro5} Dyck PJ, Thomas PK. ``Peripheral Neuropathy.'' 6th ed. Elsevier; 2021.
\bibitem{tingling_numbness:periphneuro6} Pop-Busui R, Boulton AJM, Feldman EL, et al. ``Diabetic neuropathy: a position statement by the ADA.'' Diabetes Care. 2017;40(1):136-154.
\end{thebibliography}
